\chapter{Übersicht der wichtigsten Formeln}

\section*{Einfache Zinsrechnung}
\[ Z_t = K \cdot \frac{p}{100} \cdot t = K \cdot i \cdot t \qquad Z_T = K \cdot \frac{p}{100} \cdot \frac{T}{365} \]

\section*{Zinseszinsrechnung}
\textbf{Grundformel:} \[ K_n = K_0 \cdot q^n \]
\textbf{Effektiver Zinssatz bei unterjähriger Verzinsung:} \[ i_E = (1 + i_R)^m - 1 \]

\section*{Rentenrechnung}
\textbf{Übereinstimmende Renten- und Zinsperiode:}
\[ R_n^{\text{nach}} = r \cdot \frac{q^n - 1}{q - 1} \qquad R_n^{\text{vor}} = r \cdot q \cdot \frac{q^n - 1}{q - 1} \]
\textbf{Unterjährige Renten bei jährlicher Verzinsung:}
\[ r_E = r \left(m + \tfrac{i}{2}(m - 1)\right) \qquad \bar{r}_E = r \left(m + \tfrac{i}{2}(m + 1)\right) \]

\section*{Tilgungsrechnung}
\textbf{Annuitätentilgung:} \[ S \cdot q^n = A \cdot \frac{q^n - 1}{q - 1} \]
\textbf{Ratentilgung:} Die Formeln werden in dieser Lerneinheit nicht behandelt.\\
Vorgehensweise: Festlegung einer jährlichen Tilgung, Annuität ist Summe des Tilgungsbetrages und der angefallenen Zinsen -- Aufstellen des Tilgungsplans.

\section*{Investitionsrechnung}
\textbf{Kapitalwertmethode:}
\[ C = P_1 \cdot q^{-1} + P_2 \cdot q^{-2} + \ldots + P_n \cdot q^{-n} - A_0 = \sum_{k=1}^{n} P_k \cdot q^{-k} - A_0 \]
\textbf{Methode des internen Zinsfußes:}
\[ 0 = P_1 \cdot q^{-1} + P_2 \cdot q^{-2} + \ldots + P_n \cdot q^{-n} - A_0 \]
